\documentclass[10pt,a4paper]{amsart}
\usepackage[margin=19mm]{geometry}
\usepackage{amsmath,amssymb,mathtools}
\usepackage[scr=rsfs,cal=euler]{mathalfa}
\usepackage{xfrac}
\usepackage[unicode,hidelinks,bookmarks=false]{hyperref}
\newcommand{\ie}{i.e.\ }
\newcommand{\cf}{cf.\ }
\newcommand{\resp}{respectively\ }
\newcommand{\Subsection}{\subsection}
\providecommand{\nobreakdash}{\nobreak\hbox{-}}
\setlength{\emergencystretch}{2em}
\multlinegap0pt
\usepackage{bm}
\usepackage{bbm}
\usepackage{dsfont}
\usepackage{xspace}
\usepackage{cancel}
\usepackage{mathtools}
\usepackage{amsthm}
\makeatletter
\@ifundefined{thm}{\newtheorem{thm}{Thm}}{}
\@ifundefined{prop}{\newtheorem{prop}[thm]{Proposition}}{}
\makeatother
\DeclareMathOperator{\loc}{loc}
\DeclareMathOperator{\ord}{ord}
\newcommand{\Fcal}{\mathcal{F}}
\newcommand{\Q}{\mathbb{Q}}
\newcommand{\Z}{\mathbb{Z}}
\newcommand{\cL}{\mathcal{L}}
\newcommand{\cN}{\mathcal{N}}
\newcommand{\fm}{\mathfrak{m}}
\newcommand{\ind}{\operatorname{ind}}
\newcommand{\rH}{\mathrm{H}}
\newcommand{\rank}{\mathrm{rank}}
\title[Non-vanishing of Kolyvagin systems and Iwasawa theory]{Non-vanishing of Kolyvagin systems and Iwasawa theory}
\author{Ashay Burungale, Francesc Castella, Giada Grossi,, Christopher Skinner}
\date{Source: arXiv:2312.09301v2 (CC BY 4.0)}
\begin{document}
\maketitle

\noindent\emph{Source and licence.} Non-vanishing of Kolyvagin systems and Iwasawa theory. Version arXiv:2312.09301v2 (CC BY 4.0), distributed under the Creative Commons Attribution 4.0 International licence, as recorded in the arXiv licence metadata for this exact version and shown on the abstract page https://arxiv.org/abs/2312.09301v2. Full rights notice: licenses/arxiv-2312.09301-CC-BY-4.0.txt.\par\smallskip

\noindent\textit{Editorial scope.} Counted unit: the Prop whose source label is \texttt{prop:expbound} in the article (source0.tex lines 3067--3069, source label \texttt{prop:expbound}), delivered together with the article's own complete original proof of it (source0.tex lines 3070--3099, one \texttt{proof} environment of 30 lines). No mathematics is written, repaired or re-derived: statement and proof are the article's own text line for line. Required context retained uncounted: eq:finsingkato (lines 2297--2299); prop:prime2 (lines 2990--2996). Every definition, display and label of those lines is retained unchanged. Omitted from this package and not redistributed: 1--2296, 2300--2989, 2997--3066, 3100--3203. Nothing inside a retained excerpt is omitted or altered: every retained body is delivered exactly as the article's lines are written, so no omission receipt of any kind accompanies this package. Citations: the retained excerpts carry no citation command at all, so no bibliography entry of the article is transcribed and no reference is redistributed. Reference adaptation: none was needed and none was made; every reference inside the delivered unit resolves inside it, so no unresolved reference or citation is produced. Printed slips kept literal: no printed word, symbol, hypothesis or number of the counted statement or of its proof was altered. Layout and class: the article's own class is \texttt{amsart} with packages including inputenc, amsmath, amsfonts, amssymb, graphicx, geometry, amsthm, lipsum, stmaryrd, dsfont, comment, xy, tikz-cd, hyperref; the wrapper is the lane's pinned \texttt{amsart} editorial wrapper at 10pt on A4, which restates the theorem-like environments the retained text needs and adds the article's own macro definitions that the retained text needs (\texttt{\textbackslash Fcal}, \texttt{\textbackslash Q}, \texttt{\textbackslash Z}, \texttt{\textbackslash cL}, \texttt{\textbackslash cN}, \texttt{\textbackslash fm}, \texttt{\textbackslash ind}, \texttt{\textbackslash loc}, \texttt{\textbackslash ord}, \texttt{\textbackslash rH}, \texttt{\textbackslash rank}), with extra wrapper packages (bm, bbm, dsfont, xspace, cancel, mathtools). The delivered numbering is the unit's own. Assets: no retained span contains a figure environment, a graphics-inclusion command, a PDF-page inclusion, a table or a TikZ picture; no image file is copied into this package, the package declares no asset dependency and no image file is redistributed with it. Licence: version arXiv:2312.09301v2 of this article is distributed under the Creative Commons Attribution 4.0 International licence, as recorded in the arXiv licence metadata for this exact version and shown on the abstract page https://arxiv.org/abs/2312.09301v2. A full rights notice is delivered at licenses/arxiv-2312.09301-CC-BY-4.0.txt. Characters: every retained excerpt is pure ASCII, so the delivered engine, XeLaTeX with its default Latin Modern text font, reports no missing character.
\begin{equation}\label{eq:finsingkato}
\ord(\loc_\ell(\kappa_n^{\rm Kato}))=\ord(\loc_\ell(\kappa^{\rm Kato}_{n\ell})) \ \ \text{ for every $n\in \cN, \ell\in\cL$ such that $\ell\nmid n$}.
\end{equation} 

\begin{prop}\label{prop:prime2} 
Let $k>e$ and consider two classes $c_0\in \rH^1(\Q,T^{(k)})$ and $c_1\in \rH^1(\Q,(T^{(k)})^*)$. 
Then there exist infinitely many primes $\ell\in \cL^{(k)}$ such that 
$$
\ord(\loc_\ell(c_i)) \geqslant \ord(c_i) - e, \ \ i=0,1.
$$
\end{prop}

\begin{prop}\label{prop:expbound}
Assume $\rank_{\Z_p}(\rH^1_{\Fcal}(\Q,T))=1$ and let $s_1=\ind(\kappa_1, \rH^1_{\Fcal}(\Q,T))$. For $k\gg 0$ chosen so that $k > m+s_1+2e$, we have $$\fm^{s_1+2e}\rH^1_{\Fcal^*}(\Q,(T^{(k)})^*)=0.$$
\end{prop}

\begin{proof}
Choose $k\gg s_1+2e$ such that the image of $\kappa_1$ in $\rH^1_{\Fcal}(\Q,T^{(k)})$ is non-zero and has index $s_1$. Then there exists $c_0\in\rH^1_{\Fcal}(\Q,T^{(k)})$ of order exactly $k$ such that $\varpi^{s_1}c_0 = \kappa_1$. 
%Since $\bar{T}^{G_\Q}=0$, $\rH^1_{\Fcal}(\Q,T )$ is torsion free and in particular $\rH^1_{\Fcal^*}(\Q,(T^{(k)})^*)$ has exponent strictly smaller than $k$ and we can assume $c_0$ is a generator of $\rH^1_{\Fcal}(\Q,T^{(k)})/\rH^1_{\Fcal^*}(\Q,(T^{(k)})^*)$. 
Let us write
\[
\rH^1_{\Fcal^*}(\Q,(T^{(k)})^*)=\bigoplus_{i=1}^{s} R/\varpi^{d_i}\cdot c_i, \ \ d_1\geq d_2 \geq \cdots \geq d_s.
\]
This is a cyclic-module decomposition with factors of the indicated lengths $d_i$. 

We apply Proposition \ref{prop:prime2} to the classes $c_0, c_1$ to find a prime $\ell \in \cL^{(k)}$ such that
\[
\ord(\loc_{\ell}(c_0))\geq k-e, \ \ \ \ord(\loc_{\ell}(c_1))\geq d_1-e.
\]
Recall that $\rH^1_f(\Q_{\ell},T^{(k)})$ and $\rH^1_f(\Q_{\ell}, (T^{(k)})^*)$ are isomorphic to $R/\fm^k\oplus R/\fm^t$. If $d_1<2e$, then the statement holds trivially. We assume $d_1>2e>t$. Then we have
\[
0\to \rH^1_{\Fcal^*_{\ell}}(\Q,(T^{(k)})^*) \to \rH^1_{\Fcal^*}(\Q,(T^{(k)})^*) \xrightarrow{\loc_\ell} R/\fm^{x'}\oplus R/\fm^a \to 0, \ \ \ x'\geq d_1-e, a\leq t,
\]
\[
0\to \rH^1_{\Fcal}(\Q,T^{(k)}) \to \rH^1_{\Fcal^{\ell}}(\Q,T^{(k)}) \xrightarrow{\loc_\ell} R/\fm^{k-x} \oplus R/\fm^{a'}\to 0, \ \ \ x\geq x', a'\leq t,
\]
where the second exact sequence is obtained from the first one by global duality. 

Recall that $\kappa_\ell \in \rH^1_{\Fcal^{\ell}}(\Q,T^{(k)})$ and that we cannot have $\ord(\loc_{\ell}(\kappa_\ell))\leq a'$, 
since this would contradict the assumption $k > m+s_1+2e$. Then by the Kolyvagin system relation \eqref{eq:finsingkato} and the assumptions above, we obtain
\[
k-(d_1 -e)\geq k-x\geq \ord(\loc_{\ell}(\kappa_\ell))=\ord(\loc_{\ell}(\kappa_1))=\ord(\loc_{\ell}(\varpi^{s_1}c_0))\geq k-s_1-e,
\]
which proves $\operatorname{exp}(\rH^1_{\Fcal^*}(\Q,(T^{(k)})^*))=d_1 \leq s_1+2e$. 
%Note that, if $c_0=c_1$, we get $k< s_1+2e$, contradicting the choice of $k\gg 0$.
\end{proof}

\end{document}
